% GENERATED FILE. Edit canonical lessons, then run npm run generate:books.
% canonical-source-hash: fnv1a64:d320a02668cbe75c
% canonical-lessons: ML-C381-iruttu, ML-C381-sabdam, ML-C381-nissabdata, ML-C381-mantrikkuka, ML-C381-apratyaksamakuka

\chapter{In the Dark}
\label{ch:ml-in-the-dark}
\hlchaptermodality{381}

\begin{chapteropening}
\textbf{By the end of this chapter:} \emph{I can say darkness, a sound, silence, to whisper and to disappear.}

Complete the last lesson of chapter 381: I can say darkness, a sound, silence, to whisper and to disappear.
\end{chapteropening}

\section[\texorpdfstring{\ml{ഇരുട്ട്} (iruṭṭŭ)}{ഇരുട്ട് (iruṭṭŭ)}]{\ml{ഇരുട്ട്} (iruṭṭŭ) — darkness}

\label{lesson:ML-C381-iruttu}



\begin{quote}
Before the new one: say the Malayalam for to scream, then the Malayalam for to turn back.
\end{quote}

\subsection*{You'll want to know: \ml{ഇരുട്ട്}}
\textbf{\ml{ഇരുട്ട്}} — \emph{iruṭṭŭ} — \textquotedblleft{}darkness\textquotedblright{}.

\begin{quote}
\ml{കറന്റ്} \ml{പോയി}, \ml{വീട്} \ml{മുഴുവൻ} \ml{ഇരുട്ടിലായി}.
\end{quote}

\begin{quote}
\emph{kaṟanṟŭ pōyi, vīṭŭ muḻuvan iruṭṭilāyi.}
\end{quote}

\begin{quote}
The power went and the whole house was in darkness.
\end{quote}

In the dark is \ml{ഇരുട്ടിൽ}. The adjective \ml{ഇരുണ്ട}, dark, has the same root.

\begin{cousinweb}
\textbf{\ml{ഇരുട്ട്}} (\emph{iruṭṭŭ}) is native Dravidian: DEDR 2552 lists Malayalam \emph{iruṭṭu}, \textquotedblleft{}darkness\textquotedblright{}, with \emph{iruḷ} and the verb for growing dark.
\end{cousinweb}

\begin{grammarlens}[title={What you've built}]
One more word for telling what happened.
\end{grammarlens}

\subsection*{Your turn}
\begin{itemize}
\raggedright
  \item \emph{Say it:} \emph{iruṭṭŭ}
  \item \emph{Say it:} \emph{iruṭṭŭ}, once more
  \item \emph{Say it:} \emph{kaṟanṟŭ pōyi, vīṭŭ muḻuvan iruṭṭilāyi.} — The power went and the whole house was in darkness.
  \item \emph{Recall:} say the Malayalam for to scream, then the Malayalam for to turn back, then say \emph{iruṭṭŭ} again
\end{itemize}

\begin{tcolorbox}[breakable,colback=teal!4,colframe=teal!35!black,title={Before you move on}]
What does \ml{ഇരുട്ട്} mean? (\textquotedblleft{}Darkness\textquotedblright{}.) Say it once more.
\end{tcolorbox}
\section[\texorpdfstring{\ml{ശബ്ദം} (śabdaṁ)}{ശബ്ദം (śabdaṁ)}]{\ml{ശബ്ദം} (śabdaṁ) — a sound; a voice}

\label{lesson:ML-C381-sabdam}



\begin{quote}
Before the new one: say the Malayalam for to turn back, then the Malayalam for darkness.
\end{quote}

\subsection*{You'll want to know: \ml{ശബ്ദം}}
\textbf{\ml{ശബ്ദം}} — \emph{śabdaṁ} — \textquotedblleft{}a sound; a voice\textquotedblright{}.

\begin{quote}
\ml{രാത്രി} \ml{അടുക്കളയിൽ} \ml{ഒരു} \ml{ശബ്ദം} \ml{കേട്ടു}.
\end{quote}

\begin{quote}
\emph{rātri aṭukkaḷayil oru śabdaṁ kēṭṭu.}
\end{quote}

\begin{quote}
At night a sound was heard in the kitchen.
\end{quote}

Also a voice: \ml{അവളുടെ} \ml{ശബ്ദം}, her voice.

\begin{cousinweb}
\textbf{\ml{ശബ്ദം}} (\emph{śabdaṁ}) is borrowed from Sanskrit \emph{śabda}, \textquotedblleft{}an articulate sound, a noise\textquotedblright{} (CDIAL 12298).
\end{cousinweb}

\begin{grammarlens}[title={What you've built}]
One more word for telling what happened.
\end{grammarlens}

\subsection*{Your turn}
\begin{itemize}
\raggedright
  \item \emph{Say it:} \emph{śabdaṁ}
  \item \emph{Say it:} \emph{śabdaṁ}, once more
  \item \emph{Say it:} \emph{rātri aṭukkaḷayil oru śabdaṁ kēṭṭu.} — At night a sound was heard in the kitchen.
  \item \emph{Recall:} say the Malayalam for to turn back, then the Malayalam for darkness, then say \emph{śabdaṁ} again
\end{itemize}

\begin{tcolorbox}[breakable,colback=teal!4,colframe=teal!35!black,title={Before you move on}]
What does \ml{ശബ്ദം} mean? (\textquotedblleft{}A sound; a voice\textquotedblright{}.) Say it once more.
\end{tcolorbox}
\section[\texorpdfstring{\ml{നിശ്ശബ്ദത} (niśśabdata)}{നിശ്ശബ്ദത (niśśabdata)}]{\ml{നിശ്ശബ്ദത} (niśśabdata) — silence}

\label{lesson:ML-C381-nissabdata}



\begin{quote}
Before the new one: say the Malayalam for darkness, then the Malayalam for a sound.
\end{quote}

\subsection*{You'll want to know: \ml{നിശ്ശബ്ദത}}
\textbf{\ml{നിശ്ശബ്ദത}} — \emph{niśśabdata} — \textquotedblleft{}silence\textquotedblright{}.

\begin{quote}
\ml{അറിയിപ്പ്} \ml{കേട്ടപ്പോൾ} \ml{ക്ലാസിൽ} \ml{നിശ്ശബ്ദത} \ml{പരന്നു}.
\end{quote}

\begin{quote}
\emph{aṟiyippŭ kēṭṭappōḷ klāsil niśśabdata parannu.}
\end{quote}

\begin{quote}
When they heard the notice, silence spread through the class.
\end{quote}

The adjective is \ml{നിശ്ശബ്ദമായ}, silent; the -\ml{ത} ending makes the noun.

\begin{cousinweb}
\textbf{\ml{നിശ്ശബ്ദത}} (\emph{niśśabdata}) is built from Sanskrit \emph{niḥśabda}, \textquotedblleft{}noiseless, silent\textquotedblright{} (Monier-Williams), with -\ml{ത}, the ending that makes a noun of a state. The visarga of \emph{niḥ-} becomes a second \ml{ശ} before \ml{ശബ്ദം}, so the word is written with \ml{ശ്ശ}.
\end{cousinweb}

\begin{grammarlens}[title={What you've built}]
One more word for telling what happened.
\end{grammarlens}

\subsection*{Your turn}
\begin{itemize}
\raggedright
  \item \emph{Say it:} \emph{niśśabdata}
  \item \emph{Say it:} \emph{niśśabdata}, once more
  \item \emph{Say it:} \emph{aṟiyippŭ kēṭṭappōḷ klāsil niśśabdata parannu.} — When they heard the notice, silence spread through the class.
  \item \emph{Recall:} say the Malayalam for darkness, then the Malayalam for a sound, then say \emph{niśśabdata} again
\end{itemize}

\begin{tcolorbox}[breakable,colback=teal!4,colframe=teal!35!black,title={Before you move on}]
What does \ml{നിശ്ശബ്ദത} mean? (\textquotedblleft{}Silence\textquotedblright{}.) Say it once more.
\end{tcolorbox}
\section[\texorpdfstring{\ml{മന്ത്രിക്കുക} (mantrikkuka)}{മന്ത്രിക്കുക (mantrikkuka)}]{\ml{മന്ത്രിക്കുക} (mantrikkuka) — to whisper}

\label{lesson:ML-C381-mantrikkuka}



\begin{quote}
Before the new one: say the Malayalam for a sound, then the Malayalam for silence.
\end{quote}

\subsection*{You'll want to know: \ml{മന്ത്രിക്കുക}}
\textbf{\ml{മന്ത്രിക്കുക}} — \emph{mantrikkuka} — \textquotedblleft{}to whisper\textquotedblright{}.

\begin{quote}
\ml{അവൾ} \ml{എന്റെ} \ml{ചെവിയിൽ} \ml{എന്തോ} \ml{മന്ത്രിച്ചു}.
\end{quote}

\begin{quote}
\emph{avaḷ enṟe ceviyil entō mantriccu.}
\end{quote}

\begin{quote}
She whispered something in my ear.
\end{quote}

Past: \ml{മന്ത്രിച്ചു}, \emph{mantriccu}.

Into the ear is \ml{ചെവിയിൽ}.

\begin{cousinweb}
\textbf{\ml{മന്ത്രിക്കുക}} (\emph{mantrikkuka}) is built from Sanskrit \emph{mantra}, \textquotedblleft{}thought, counsel, a spell\textquotedblright{} (CDIAL 9834), with the Malayalam verb ending -\ml{ിക്കുക}: to say something low, as a spell is said.
\end{cousinweb}

\begin{grammarlens}[title={What you've built}]
One more word for telling what happened.
\end{grammarlens}

\subsection*{Your turn}
\begin{itemize}
\raggedright
  \item \emph{Say it:} \emph{mantrikkuka}
  \item \emph{Say it:} \emph{mantrikkuka}, once more
  \item \emph{Say it:} \emph{avaḷ enṟe ceviyil entō mantriccu.} — She whispered something in my ear.
  \item \emph{Recall:} say the Malayalam for a sound, then the Malayalam for silence, then say \emph{mantrikkuka} again
\end{itemize}

\begin{tcolorbox}[breakable,colback=teal!4,colframe=teal!35!black,title={Before you move on}]
What does \ml{മന്ത്രിക്കുക} mean? (\textquotedblleft{}To whisper\textquotedblright{}.) Say it once more.
\end{tcolorbox}
\section[\texorpdfstring{\ml{അപ്രത്യക്ഷമാകുക} (apratyakṣamākuka)}{അപ്രത്യക്ഷമാകുക (apratyakṣamākuka)}]{\ml{അപ്രത്യക്ഷമാകുക} (apratyakṣamākuka) — to disappear}

\label{lesson:ML-C381-apratyaksamakuka}



\begin{quote}
Before the new one: say the Malayalam for silence, then the Malayalam for to whisper.
\end{quote}

\subsection*{You'll want to know: \ml{അപ്രത്യക്ഷമാകുക}}
\textbf{\ml{അപ്രത്യക്ഷമാകുക}} — \emph{apratyakṣamākuka} — \textquotedblleft{}to disappear\textquotedblright{}.

\begin{quote}
\ml{കള്ളൻ} \ml{ഇരുട്ടിൽ} \ml{അപ്രത്യക്ഷമായി}.
\end{quote}

\begin{quote}
\emph{kaḷḷan iruṭṭil apratyakṣamāyi.}
\end{quote}

\begin{quote}
The thief disappeared into the darkness.
\end{quote}

Past: \ml{അപ്രത്യക്ഷമായി}, \emph{apratyakṣamāyi}.

\ml{അ}-, not, + \ml{പ്രത്യക്ഷം}, before the eyes, + \ml{ആകുക}, to become: to go out of sight.

\begin{cousinweb}
\textbf{\ml{അപ്രത്യക്ഷമാകുക}} (\emph{apratyakṣamākuka}) is built from Sanskrit \emph{pratyakṣa}, \textquotedblleft{}visible, before the eyes\textquotedblright{} (CDIAL 8635), with \emph{a-}, \textquotedblleft{}not\textquotedblright{}, and the Malayalam \ml{ആകുക}, \textquotedblleft{}to become\textquotedblright{}.
\end{cousinweb}

\begin{grammarlens}[title={What you've built}]
One more word for telling what happened.
\end{grammarlens}

\subsection*{Your turn}
\begin{itemize}
\raggedright
  \item \emph{Say it:} \emph{apratyakṣamākuka}
  \item \emph{Say it:} \emph{apratyakṣamākuka}, once more
  \item \emph{Say it:} \emph{kaḷḷan iruṭṭil apratyakṣamāyi.} — The thief disappeared into the darkness.
  \item \emph{Recall:} say the Malayalam for silence, then the Malayalam for to whisper, then say \emph{apratyakṣamākuka} again
\end{itemize}

\begin{tcolorbox}[breakable,colback=teal!4,colframe=teal!35!black,title={Before you move on}]
What does \ml{അപ്രത്യക്ഷമാകുക} mean? (\textquotedblleft{}To disappear\textquotedblright{}.) Say it once more.
\end{tcolorbox}
